% methods_embeddings.tex
\subsection{Embeddings}

The main analyses use precomputed item embeddings from Qwen3-Embedding-8B \citep{qwen-2025}, a 4,096-dimensional text-embedding model. Each of the 50 item texts was embedded once, offline, with last-token pooling, no instruction prefix, and no normalization. The resulting $50 \times 4096$ matrix (with texts, codes, domain labels, and scoring key) ships as \texttt{data(big5)}, so the analyses need no Python or network access.

Label anchoring and semantic projection also require vectors for non-item words: the five construct names (\emph{extraversion}, \emph{neuroticism}, \emph{agreeableness}, \emph{conscientiousness}, \emph{openness}) and the projection poles (\emph{anxious}, \emph{calm}, \emph{solitary}, \emph{sociable}). A script in the reproduction archive extracted these from a precomputed one-million-word lexicon embedded with the same model, placing them in the items' vector space as the cosine comparisons require.

The package's on-device default for new text is the smaller Qwen3-Embedding-0.6B \citep{qwen-2025}, yielding $50 \times 1024$ vectors for the same items, and backs the two demonstrations that must embed text live: the \texttt{sfa\_embed()} re-embedding check, which compares fresh vectors with same-model offline reference vectors loaded from a NumPy \texttt{.npz} archive by \texttt{sfa\_load\_npz()} (the package's import route for external embeddings), and \texttt{sfa\_item\_fit()}'s vetting of candidate items that postdate the archives. Both run through \texttt{reticulate} and \texttt{sentence-transformers}, which downloads the model from the Hugging Face Hub on first use.

Finally, one similarity matrix is not embedding-based: \texttt{sfa\_nli\_matrix()} scores all item pairs with \texttt{cross-encoder/nli-deberta-v3-base}, a natural-language-inference (NLI) cross-encoder trained on SNLI-tradition corpora \citep{bowman-etal-2015}. The next subsection describes its role.
